\section{Experiments and Results}
In this section, we conduct extensive experiments to evaluate the performance of the proposed Cross-Modality Semantic Alignment Transformer (CMSA-Trans) for zero-shot fault diagnosis of rotating machinery. The evaluation focuses on the model's ability to generalize to unseen fault categories by leveraging LLM-generated semantic descriptions.

\subsection{Experimental Setup}
To validate the effectiveness of the proposed method, we utilize two widely recognized mechanical fault diagnosis benchmarks: the Case Western Reserve University (CWRU) bearing dataset and the South East University (SEU) gearbox dataset. These datasets provide a comprehensive evaluation across different rotating components (bearings vs. gears) and varying levels of background noise and mechanical complexity.

The Case Western Reserve University (CWRU) bearing dataset consists of vibration measurements from a 2 hp Reliance Electric motor under four load conditions (0--3 hp), sampled at 12 kHz. Our corpus includes 3,500 signal segments labeled by fault diameter and component position. For zero-shot evaluation, seen classes include Normal (N), Ball fault (B, 0.007''), and Outer Race fault (OR, 0.014'') at 1797 rpm, while Inner Race faults (IR, 0.021'') and complex multi-component faults (0.028'' ball) are unseen categories. The seen classes are split 80/20 for training (2,000 samples) and validation (500 samples), with 1,000 unseen samples reserved for zero-shot evaluation.

The SEU gearbox dataset provides a more challenging testbed featuring a multi-stage transmission system (motor, planetary gear reducer, parallel shaft gearbox, and magnetic powder brake). Data is sampled at 5120 Hz across rotating speeds of 20 Hz and 30 Hz under loads of 0 V and 7.3 V, yielding 4,000 segments. Five health states are characterized: Normal, Chipped tooth, Worn tooth, Broken tooth, and Root fault. Normal, Broken, and Chipped states serve as seen classes (2,400 samples), while Worn and Root faults are unseen (1,600 samples). Semantic descriptors for unseen faults are generated by GPT-4 using expert-level diagnostic prompts covering concepts such as modulation patterns, circular pitch errors, and mesh frequency sidebands.

Raw vibration signals are segmented using an overlapping sliding window (size: 1024, step: 512) to preserve transient features while ensuring data diversity. Each segment undergoes Z-score normalization to eliminate DC bias and minimize the influence of varying operating conditions. A Hanning window is applied to reduce spectral leakage before feature extraction. The resulting normalized sequences $X \in \mathbb{R}^{1024 \times 1}$ serve as input tensors, yielding approximately 2,500 seen-class samples and 1,000 unseen-class samples per dataset. Seen-class samples are partitioned 80/20 for training and validation.

\subsection{Implementation Details}
The CMSA-Trans is implemented in PyTorch 2.1, comprising a 3-layer 1D-CNN backbone (kernel size 3, stride 1) followed by a 4-layer Transformer encoder with 8-head self-attention. The embedding dimension $d_{model}$ is set to 256, balancing representational capacity and computational efficiency.

For the language branch, pre-trained BERT-large produces 1024-dimensional semantic embeddings, which are projected into the shared 256-dimensional space via a two-layer FFN. The model is optimized using Adam (weight decay $1e{-4}$) with an initial learning rate of $0.001$ and Cosine Annealing scheduling over 200 epochs. Batch size is 64. Training is conducted on an NVIDIA RTX 3090 GPU, requiring approximately 55 minutes per dataset. The trade-off hyperparameter $\lambda$ in Eq.9 is set to 0.5. All experiments are repeated 10 times, and we report mean accuracy with standard deviation.


\subsection{Zero-Shot Recognition Performance}
The zero-shot diagnostic performance of the CMSA-Trans is rigorously benchmarked against five state-of-the-art (SOTA) baseline methods to highlight its superiority in cross-modality alignment. These baselines include: (1) WDCNN \cite{zhang2017new}, a classic CNN-based model adapted for ZSL via attribute matching; (2) ZSL-GAN \cite{xian2019fvaegan}, a generative adversarial network designed to synthesize unseen class features to bridge the domain gap between observed and novel faults; (3) Deep Multi-modal ZSL (DMZSL) \cite{schonfeld2019generalized}, which uses a shared subspace for vision and language to facilitate knowledge transfer through common feature representations; (4) Attribute-based ZSL (AZSL) \cite{lampert2014attribute}, relying on manual engineer-defined fault attributes such as frequency bands and impact intensities; and (5) Bi-LSTM-ZSL \cite{chen2019bearing}, which captures temporal dependencies but lacks specialized semantic alignment layers for cross-modal interaction. As summarized in Table~\ref{tab:comparison}, CMSA-Trans demonstrates a significant performance advantage across all fault categories \cite{li2022zero, lu2017deep}.

On the CWRU bearing dataset, CMSA-Trans achieves an impressive average zero-shot accuracy of 89.4\% for the unseen fault categories. This represents a 4.2% improvement over the runner-up method (ZSL-GAN). While GAN-based methods struggle with the instability of the min-max game and often generate noisy feature distributions that lack physical consistency \cite{lei2020applications}, the Transformer-based CMSA-Trans architecture provides a stable and structured mapping between the vibration signal patches and the LLM semantic tokens. For the "Inner Race Fault" (unseen), CMSA-Trans achieves 91.2% accuracy, significantly outperforming AZSL (76.8%) and WDCNN (71.3%). This improvement is due to the LLM's ability to describe the "higher-order harmonics" and "envelope modulation" typically seen in inner race defects, which are often overlooked in simplified manual attribute sets. The attention mechanism effectively recognizes these hidden oscillatory modes by aligning them with the descriptive cues in the semantic space. Deep architectures like WDCNN often fail to capture these subtle signatures because their fixed kernels are insensitive to the long-range temporal periodicity of inner race impacts \cite{hochreiter1997long}. In contrast, our model treats each impact as a distinct semantic token, allowing for a more granular identification process.

On the more complex SEU gearbox dataset, the advantages of our semantic alignment strategy are even more pronounced. The SEU dataset involves compound vibrations from multiple gears and bearings within the gearbox enclosure, making the distinction between "worn" and "root" faults extremely difficult without domain-specific expert knowledge. CMSA-Trans achieves a zero-shot accuracy of 82.5% on this dataset, whereas traditional attribute-based methods (AZSL) drop to 65.4%. The performance gap signifies that the rich, descriptive technical semantics provided by LLMs act as a superior bridge for knowledge transfer compared to sparse binary attributes. Specifically, for the "Broken Tooth" category, CMSA-Trans maintains 88.7% accuracy, which is 15.2% higher than Bi-LSTM-ZSL. This is because root faults exhibit unique spectral signatures in the sideband frequencies that are explicitly detailed in our LLM prompts but are ignored by LSTM-based sequence models that only look at global temporal patterns \cite{zhang2021attention}. Fig.~\ref{fig:accuracy_bar} presents the confusion matrix for the unseen categories, showing that the model rarely confuses the broken tooth fault with the normal state, thanks to the distinct semantic prototypes provided by the language model. The high precision and recall across unseen categories demonstrate that CMSA-Trans is highly effective at generalizing to novel mechanical health states without requiring any labeled signal data during training, showcasing its potential for practical prognostic and health management (PHM) applications \cite{jia2016deep}.

\subsection{Ablation Study}
To further dissect the contribution of each component within the CMSA-Trans framework, we conduct a systematic ablation study on both the CWRU and SEU datasets. We evaluate four distinct model configurations: (1) CMSA-Trans-Full: the proposed architecture with all components active; (2) CMSA-Trans-NoLLM: utilizes traditional 20-dimensional manual attribute vectors instead of 1024-dimensional LLM-generated semantic embeddings; (3) CMSA-Trans-NoAttn: replaces the cross-modality attention alignment module with a simple global average pooling followed by a concatenation fusion of signal and text features; (4) CMSA-Trans-Base: replaces the Transformer encoder with a standard 1D-ResNet structure, thereby removing the self-attention mechanism for temporal dependency modeling.

The quantitative results of this study across both datasets highlight a consistent trend. For the CWRU dataset, CMSA-Trans-Full achieves 89.4% accuracy, while CMSA-Trans-NoLLM, CMSA-Trans-NoAttn, and CMSA-Trans-Base achieve 81.6%, 83.9%, and 85.1%, respectively. On the more challenging SEU dataset, CMSA-Trans-Full maintains a high 82.5%, whereas the variants drop significantly to 74.7%, 77.0%, and 78.4%. As shown in Table~\ref{tab:ablation}, the removal of LLM-generated semantics (CMSA-Trans-NoLLM) results in the most significant performance degradation, with a 7.8% drop in zero-shot accuracy on the SEU dataset. This confirms our hypothesis that the "semantic bottleneck" in traditional ZSL for fault diagnosis is caused by the low information density and subjective nature of manual attributes. LLMs provide a high-dimensional, continuous semantic space that better preserves the transitive relations between different fault types by encoding linguistic relationships. For instance, the semantic distance between "worn tooth" and "chipped tooth" is much more accurately captured in the 1024-dimensional BERT embedding than in a binary attribute vector where both might share 90% of the same labels.

Furthermore, the exclusion of the cross-attention mechanism (CMSA-Trans-NoAttn) leads to a 5.5% accuracy loss. This indicates that local alignment between specific signal segments (e.g., transient impacts) and semantic keywords is vital for fine-grained fault recognition. Simple global fusion fails to capture the intricate temporal-semantic correlations. Without the attention layer, the model attempts to map the entire averaged signal vector to the entire semantic prototype, losing the ability to associate specific periodic pulses with specific textual descriptors of those pulses. Finally, the comparison with CMSA-Trans-Base shows that the Transformer's self-attention layers are better at capturing the long-range periodic nature of rotating machinery vibrations compared to the local receptive fields of standard residual convolutions, which tend to focus on stationary features. The ResNet-based backbone (CMSA-Trans-Base) struggles with signals collected under varying load conditions, where the fault impacts vary in frequency and phase. The Transformer's ability to attend to all time steps simultaneously allows it to remain invariant to these phase shifts, which is essential for consistent zero-shot alignment across different operating speeds. This analysis demonstrates that the synergy between technical language model semantics and global temporal attention is the core driver of our framework's success in industrial zero-shot scenarios.

\subsection{Parameter Sensitivity & Robustness}
We conduct a sensitivity analysis to determine the optimal configuration of the latent embedding space and evaluate the model's robustness against measurement noise. Fig.~\ref{fig:sensitivity} illustrates the variation in zero-shot accuracy as the hidden embedding dimension $d_{model}$ is tuned from 64 to 512. We observe that increasing the dimension from 64 to 256 yields a sharp improvement in performance as the model gains the capacity to represent complex cross-modal relationships between the vibration and semantic domains. However, a further increase to 512 results in a slight decrease in validation accuracy, likely due to the increased dimensionality making the projection head prone to overfitting the limited training samples in the seen classes. Consequently, 256 is selected as the optimal dimension for our final architecture to ensure both generalization and stability. We further investigate the impact of the training batch size on convergence behavior and final accuracy. With a batch size of 32, the model exhibits high gradient variance, leading to oscillatory convergence and a 2.1\% accuracy reduction on the SEU dataset. Increasing the batch size to 128 stabilizes training but slightly degrades generalization due to the reduced stochasticity in gradient updates, resulting in a 0.8\% accuracy drop compared to the selected batch size of 64. Additionally, we evaluate the sensitivity to the learning rate by testing values from $5 \times 10^{-4}$ to $5 \times 10^{-3}$. A learning rate of $5 \times 10^{-3}$ causes divergence in the cross-modal alignment loss during the first 20 epochs, while $5 \times 10^{-4}$ leads to premature convergence at a suboptimal local minimum with 3.4\% lower accuracy. The selected learning rate of $10^{-3}$ with cosine annealing provides the best trade-off between convergence speed and final performance across both datasets.

The robustness of the proposed CMSA-Trans is a critical factor for its deployment in harsh industrial environments. To test this, we inject Additive White Gaussian Noise (AWGN) into the raw test signals to simulate different Signal-to-Noise Ratios (SNR), ranging from -10 dB to 10 dB. As shown in the robustness curve in Fig.~\ref{fig:snr_curve}, CMSA-Trans maintains a relatively stable performance even in low SNR conditions. At SNR = 0 dB, CMSA-Trans achieves 75.6% accuracy, which is 12.3% higher than the performance of WDCNN under the same conditions. This superior noise resistance is attributed to two factors: the multi-scale CNN backbone that extracts noise-invariant features at different temporal resolutions, and the Transformer-based attention mechanism that acts as a temporal filter, suppressing random noise while focusing on the quasi-periodic fault impacts. These results emphasize that CMSA-Trans is not only accurate in ideal laboratory conditions but also highly reliable for actual field diagnosis where noise interference is inevitable, often obfuscating the diagnostic signatures of incipient faults.

\subsection{Visual Analysis}
We provide visual interpretation of the alignment mechanism using t-SNE projections and attention heatmaps. As illustrated in Fig.~\ref{fig:tsne}, the 256-dimensional features from the alignment layer are projected into 2D space, where unseen class signal features (Inner Race and Combined Faults) form distinct clusters spatially adjacent to their corresponding LLM-generated semantic prototypes. The most common misclassification occurs between "Worn Tooth" and "Root Fault" categories due to overlapping low-frequency sideband characteristics. However, CMSA-Trans reduces this inter-class confusion rate to 6.3\%, compared to 18.9\% for AZSL, by leveraging discriminative semantic descriptions that differentiate "gradual surface degradation" from "stress concentration at the tooth root." The spatial distances between fault clusters in the embedding space mirror the semantic distances between their textual descriptions, validating that the model learns a physically meaningful semantic manifold.

The attention heatmaps (Fig.~\ref{fig:attn_map}) further reveal that when presented with an unseen inner race fault signal, the attention heads consistently assign high weights to semantic tokens such as "periodic peak," "carrier frequency," and "bearing housing resonance." This demonstrates that CMSA-Trans establishes deep logical connections between physical failure characteristics and natural language descriptions, providing a transparent and interpretable foundation for autonomous fault diagnosis in smart manufacturing systems.

